\exercisesfor{5}

In the following exercises we assume the hypotheses and notation of §5. F denotes the free group \(\mathrm{F}(\mathbf{X})\) and g the unique homomorphism of F into the Magnus group \(\Gamma(\mathbf{X})\) such that \(g(x)=1+x\) for all \(x\in X\) (cf.~Theorem 1).

\begin{enumerate}
\def\labelenumi{\arabic{enumi}.}
\tightlist
\item
  Let the algebra \(\mathbf{K}[\mathbf{F}]\) of \(\mathbf{F}\) be given the filtration \((\mathbf{I}^n)\) consisting of the powers of the augmentation ideal \(\mathbf{I}\) ( \(\S 4\) , Exercise 6).
\end{enumerate}

\begin{enumerate}
\def\labelenumi{(\alph{enumi})}
\item
  Let \(\tilde{g}\colon \mathrm{K}[\mathrm{F}] \to \tilde{\mathrm{A}}(\mathrm{X})\) be the unique algebra homomorphism extending \(g\colon \mathrm{F} \to \tilde{\mathrm{A}}(\mathrm{X})^*\) . Show that \(\tilde{g}\) maps \(\mathrm{I}^{\mathfrak{n}}\) to the ideal \(\tilde{\mathrm{A}}_{\mathfrak{n}}(\mathrm{X})\) (cf.~no. 1) and defines when taking quotients an isomorphic \(\tilde{g}_{n}\) of \(\mathrm{K}[\mathrm{F}]/\mathrm{I}^{\mathfrak{n}}\) onto \(\tilde{\mathrm{A}}(\mathrm{X})/\tilde{\mathrm{A}}_{\mathfrak{n}}(\mathrm{X})\) . (Define an inverse homomorphism of \(\tilde{g}_{n}\) by means of the homomorphism of \(\mathrm{A}(\mathrm{X})\) into \(\mathrm{K}[\mathrm{F}]\) which maps \(x\) to \(x - 1\) for all \(x \in \mathrm{X}\) .) Deduce that \(\tilde{\mathrm{A}}(\mathrm{X})\) is isomorphic to the Hausdorff completion of \(\mathrm{K}[\mathrm{F}]\) under the topology defined by \((\mathrm{I}^{\mathfrak{n}})\) .
\item
  Suppose that \(\mathbf{K} = \mathbf{Z}\) . Show that the filtration \((\mathrm{I}^{\mathrm{n}})\) is separated (use Proposition 3 and Exercise 6 (e) of §4) and that the filtration of \(\mathbf{F}\) defined by it (Exercise 6 (a) of §4) coincides with the filtration \((\mathbf{C}^{\mathrm{n}}\mathbf{F})\) . Deduce that \(\tilde{g}\colon \mathbf{Z}[\mathrm{F}] \to \hat{\mathrm{A}}_{\mathbf{Z}}(\mathrm{X})\) is injective.
\end{enumerate}

\begin{enumerate}
\def\labelenumi{\arabic{enumi}.}
\setcounter{enumi}{1}
\tightlist
\item
  Let G be a group, let K{[}G{]} be its algebra over K and let I be its augmentation ideal (§ 4, Exercise 6). Let \(\varepsilon\) denote the canonical homomorphism of K{[}G{]} onto K; then \(\operatorname{Ker}(\varepsilon) = I\) .
\end{enumerate}

\begin{enumerate}
\def\labelenumi{(\alph{enumi})}
\tightlist
\item
  Let M be a left K{[}G{]}-module and let Z(G, M) be the group of crossed homomorphisms of G into M (Algebra, Chapter I, § 6, Exercise 7). If \(f \in \operatorname{Hom}_{\mathbb{K}[\mathbb{G}]}(\mathrm{I}, \mathrm{M})\) , let \(f_1\) be the mapping \(g \mapsto f(g - 1)\) of \(\mathbf{G}\) into \(\mathbf{M}\) . Show that \(f_1\) is a crossed homomorphism (use the identity
\end{enumerate}

\[
g g ^ {\prime} - 1 = g (g ^ {\prime} - 1) + g - 1
\]

to show that \(f_{1}(gg') = g \cdot f_{1}(g') + f_{1}(g)\) and that \(f \mapsto f_{1}\) is an isomorphism of \(\operatorname{Hom}_{\mathbb{K}[G]}(\mathbf{I}, \mathbf{M})\) onto \(\mathbf{Z}(\mathbf{G}, \mathbf{M})\) .

\begin{enumerate}
\def\labelenumi{(\alph{enumi})}
\setcounter{enumi}{1}
\item
  Take G to be the free group \(\mathrm{F} = \mathrm{F}(\mathrm{X})\) . Show that, for every mapping \(\theta: \mathrm{X} \to \mathrm{M}\) , there exists one and only one element \(f_{\theta}\) of \(\mathrm{Z}(\mathrm{M}, \mathrm{F})\) which extends \(\theta\) . (Use the interpretation of \(\mathrm{Z}(\mathrm{F}, \mathrm{M})\) in terms of the semi-direct product of F by M, cf.~Algebra, Chapter I, loc. cit.) Deduce, using (a), that the family \((x - 1)_{x \in \mathbb{X}}\) is a basis of I as a left K{[}F{]}-module.
\item
  By (b) every element \(u \in \mathbf{K}[\mathbf{F}]\) can be written uniquely in the form
\end{enumerate}

\[
u = \varepsilon (u) + \sum_ {x \in \mathbf {X}} \mathrm{D} _ {x} (u) (x - 1),
\]

where \(\mathrm{D}_{x}(u)\in\mathrm{K}[\mathrm{F}]\) and the \(\mathrm{D}_{x}(u)\) are zero except for a finite number. The mapping \(u\mapsto\mathrm{D}_{x}(u)\) is called the partial derivative with respect to x. It is characterized by the following properties:

\begin{enumerate}
\def\labelenumi{(\roman{enumi})}
\item
  \(D_{x}\) is a K-linear mapping of K{[}F{]} into K{[}F{]}.
\item
  \(\mathrm{D}_{x}(uv)=u.\mathrm{D}_{x}(v)+\varepsilon(v).\mathrm{D}_{x}(u)\) for u, v in K{[}F{]}.
\item
  \(\mathrm{D}_{x}(x)=1\) and \(\mathrm{D}_{x}(y)=0\) if \(y\in X-\{x\}\) .
\end{enumerate}

If \(n\geqslant 1\) , then

\[
\begin{array}{c} \mathrm{D} _ {x} (x ^ {n}) = 1 + x + \dots + x ^ {n - 1} \\ \mathrm{D} _ {x} (x ^ {- n}) = - x ^ {- n} - x ^ {- n + 1} - \dots - x ^ {- 1}. \end{array}
\]

\begin{enumerate}
\def\labelenumi{(\alph{enumi})}
\setcounter{enumi}{3}
\tightlist
\item
  Let \(\varepsilon_{\mathrm{A}}\) be the canonical homomorphism of A(X) onto K. Show that every element \(u\) of A(X) can be written uniquely in the form
\end{enumerate}

\[
u = \varepsilon_ {\mathrm{A}} (u) + \sum_ {x \in \mathbf {X}} d _ {x} (u). x,
\]

where \(d_{x}(u) \in \mathrm{A}(\mathrm{X})\) and the \(d_{x}(u)\) are zero except for a finite number. Show that \(d_{x}: u \to d_{x}(u)\) is an endomorphism of \(\mathrm{A}(\mathrm{X})\) which can be extended by continuity to \(\hat{\mathrm{A}}(\mathrm{X})\) and enjoys properties analogous to (i), (ii), (iii) above. Show \(d_{x} \circ \tilde{g} = \tilde{g} \circ D_{x}\) , where \(\tilde{g}\) is the homomorphism of K{[}F{]} into \(\hat{\mathrm{A}}(\mathrm{X})\) which extends g (cf.~Exercise 1).†

¶ 3. We preserve the notation of Exercise 2.

\begin{enumerate}
\def\labelenumi{(\alph{enumi})}
\tightlist
\item
  Let \(J\) be a right ideal of \(K[F]\) contained in \(I\) and let \(u\) be an element of \(I\) . Show that
\end{enumerate}

\[
u \in J. I \Leftrightarrow D _ {x} (u) \in J \quad \text { for   all } x \in X.
\]

In particular, an element of I belongs to \(\mathbf{I}^n\) ( \(n \geqslant 1\) ) if and only if all its partial derivatives belong to \(\mathbf{I}^{n-1}\) .


\begin{enumerate}
\def\labelenumi{(\alph{enumi})}
\setcounter{enumi}{1}
\tightlist
\item
  Let R be a normal subgroup of F, let G = F/R and let J be the kernel of the canonical homomorphism \(\gamma: K[F] \to K[G]\) . Show that J is generated as a left (resp. right) ideal by the elements r - 1, where \(r \in R\) . Prove the exactness of the sequence
\end{enumerate}

\[
0 \longrightarrow \mathrm{J} / (\mathrm{J}. \mathrm{I}) \stackrel {\alpha} {\longrightarrow} \mathrm{I} / (\mathrm{J}. \mathrm{I}) \stackrel {\beta} {\longrightarrow} \mathrm{K} [ \mathrm{G} ] \stackrel {\delta} {\longrightarrow} \mathrm{K} \longrightarrow 0,\tag{*}
\]

where \(\alpha\) is derived by taking quotients from the inclusion of J in I and \(\beta\) is derived by taking quotients from the restriction of \(\gamma\) to I.

\begin{enumerate}
\def\labelenumi{(\alph{enumi})}
\setcounter{enumi}{2}
\item
  For \(u \in \mathbf{I}\) , \(x \in \mathbf{X}\) , let \(\overline{\mathrm{D}}_x(u)\) denote the image of \(\mathrm{D}_x(u)\) in \(\mathrm{K}[\mathrm{G}]\) under \(\gamma\) . Show, using (a), that the family \((\overline{\mathrm{D}}_x)_{x \in \mathbf{X}}\) defines on taking quotients an isomorphism of \(\mathrm{I} / (\mathrm{J}. \mathrm{I})\) onto \(\mathrm{K}[\mathrm{G}]^{(\mathbb{X})}\) .
\item
  For \(r \in \mathbb{R}\) , let \(\theta(r)\) be the image of \(r - 1\) in \(J/(J.I)\) . Show that
\end{enumerate}

\[
\theta (r r ^ {\prime}) = \theta (r) + \theta (r ^ {\prime}) \quad \text { and } \quad \theta (y r y ^ {- 1}) = y. \theta (r) \quad \text { for   } r, r ^ {\prime} \text {   in   } R, y \in F.
\]

(Use the identities

\[
\begin{array}{c} {r r ^ {\prime} - 1 = (r - 1) (r ^ {\prime} - 1) + (r - 1) + (r ^ {\prime} - 1)} \\ {y r y ^ {- 1} - 1 = y (r - 1) (y ^ {- 1} - 1) + y (r - 1).)} \end{array}
\]

Deduce that \(\theta\) defines a homomorphism of \(R/(R, R)\) into \(J/(J, I)\) compatible with the action of \(G\) and show that the image of this homomorphism generates the K-module \(J/(J, I)\) ; when \(K = Z\) , show that an isomorphism is thus obtained of \(R/(R, R)\) onto \(J/(J, I)\) (define the inverse homomorphism directly). (e) Let \((r_{\alpha})_{\alpha \in A}\) be a family of elements of \(R\) generating \(R\) as a normal subgroup of \(F\) . Show that the \(\theta(r_{\alpha})\) generate the K{[}G{]}-module \(J/(J, I)\) .

\begin{enumerate}
\def\labelenumi{(\alph{enumi})}
\setcounter{enumi}{5}
\tightlist
\item
  The matrix \((\overline{\mathrm{D}}_x(r_\alpha))_{x\in \mathbf{X},\alpha \in \mathbf{A}}\) defines a homomorphism
\end{enumerate}

\[
\rho \colon \mathrm{K} [ \mathrm{G} ] ^ {(\mathrm{A})} \rightarrow \mathrm{K} [ \mathrm{G} ] ^ {(\mathrm{X})}
\]

of left K{[}G{]}-modules. Show that the sequence

\[
\mathrm{K} [ \mathrm{G} ] ^ {(\mathrm{A})} \xrightarrow {\rho} \mathrm{K} [ \mathrm{G} ] ^ {(\mathrm{X})} \xrightarrow {\delta} \mathrm{K} [ \mathrm{G} ] \xrightarrow {\varepsilon} \mathrm{K} \longrightarrow 0\tag{**}
\]

is exact, where \(\delta\) is the homomorphism \((u_x) \mapsto \sum_{x \in X} u_x(\gamma(x) - 1)\) .

(Transform the exact sequence (*) using (c), (d), (e) above.)†

\begin{enumerate}
\def\labelenumi{\arabic{enumi}.}
\setcounter{enumi}{3}
\tightlist
\item
  For all \(n \in \mathbf{N}\) , let \(\binom{\mathrm{T}}{n}\) denote the polynomial
\end{enumerate}

\[
\mathrm{T} (\mathrm{T} - 1) \dots (\mathrm{T} - n + 1) / n!.
\]

If \(f(\mathrm{T})\) is a polynomial, let \(\Delta f\) denote the polynomial \(f(\mathrm{T} + 1) - f(\mathrm{T})\) .

Then \(\Delta\binom{\mathrm{T}}{0} = \Delta 1 = 0\) and \(\Delta\binom{\mathrm{T}}{n} = \binom{\mathrm{T}}{n - 1}\) if \(n \geqslant 1\) .

\begin{enumerate}
\def\labelenumi{(\alph{enumi})}
\item
  Let \(f \in \mathbf{Q}[\mathrm{T}]\) . Show the equivalence of the following properties:
\item
  \(f\) maps \(\mathbf{Z}\) into itself.
\end{enumerate}

\begin{enumerate}
\def\labelenumi{(\roman{enumi})}
\setcounter{enumi}{1}
\item
  \(f\) is a linear combination with coefficients in \(\mathbf{Z}\) of the \(\binom{\mathrm{T}}{n}\) .
\item
  \(f(0)\in \mathbf{Z}\) and \(\Delta f\) maps \(\mathbf{Z}\) into itself.
\end{enumerate}

(Argue by induction on \(\deg(f)\) , noting that \(\deg(\Delta f) = \deg(f) - 1\) if \(\deg(f) \neq 0\) .)

A polynomial \(f\) satisfying the above properties is called binomial. The sum, product and composition of two binomial polynomials is a binomial polynomial.

\begin{enumerate}
\def\labelenumi{(\alph{enumi})}
\setcounter{enumi}{1}
\item
  Let \(p\) be a prime number and let \(f\) be a binomial polynomial. Show that \(f\) maps the ring \(\mathbf{Z}_p\) of \(p\) -adic integers into itself (use the continuity of \(f\) and the fact that \(\mathbf{Z}\) is dense in \(\mathbf{Z}_p\) ).
\item
  Let \(\mathbf{P}\) be a set of prime numbers, let \(\mathrm{S}_P\) be the set of P-integers (§ 4, Exercise 14) and let \(\mathbf{Z}_P = \mathrm{S}_{\mathrm{P}}^{-1}\mathbf{Z}\) . Show that if \(f\) is a binomial polynomial, \(f\) maps \(\mathbf{Z}_{\mathrm{P}}\) into itself (apply (b) to the prime numbers \(p\) not belonging to \(\mathbf{P}\) ).
\end{enumerate}

\begin{enumerate}
\def\labelenumi{\arabic{enumi}.}
\setcounter{enumi}{4}
\tightlist
\item
  Let \(\Gamma = \Gamma(X)\) be the Magnus group on \(X\) and let \((\Gamma_n)\) be its natural filtration (no. 2). Show that the graded Lie algebra \(\mathrm{gr}(\Gamma)\) corresponding to this filtration is isomorphic to the Lie subalgebra \(\bigoplus_{n\geqslant 1}A_n(X)\) of \(A(X)\) . If \(X\neq \emptyset\) , this Lie algebra is not generated by its elements of degree 1; deduce that \((\Gamma_n)\) is not the lower central series of \(\Gamma\) .
\end{enumerate}

¶ 6. Let P be a set of prime numbers, let \(S_P\) be the set of P-integers (§ 4, Exercise 14) and let \(Z_P = S_P^{-1}Z\) .

\begin{enumerate}
\def\labelenumi{(\alph{enumi})}
\tightlist
\item
  Suppose that, for all \(s \in S_{P}\) , the mapping \(k \mapsto sk\) is a bijection of K onto itself; this is equivalent to saying that K can be given a \(Z_{P}\) -algebra structure.
\end{enumerate}

With the notation of Exercise 5, show that, for all \(s \in S_{\mathbb{P}}\) , the mapping \(a \mapsto a^s\) of \(\Gamma\) into itself is bijective and that the same is true for each quotient \(\Gamma / \Gamma_n\) for \(n \geqslant 1\) . If \(t \in \mathbf{Z}_{\mathbb{P}}\) and \(a \in \Gamma\) (resp. \(a \in \Gamma / \Gamma_n\) ), \(a^s\) is defined as in Exercise 16 of §4; writing \(a\) in the form \(1 + \alpha\) , where \(\alpha \in \hat{\mathrm{A}}_1(\mathrm{X})\) , show that

\[
a ^ {t} = (1 + \alpha) ^ {t} = \sum_ {n = 0} ^ {\infty} \binom {t} {n} \alpha^ {n}.
\]

(Note that the coefficients \(\binom{t}{n}\) belong to \(\mathbf{Z}_{\mathrm{P}}\) , cf.~Exercise 4.)

\begin{enumerate}
\def\labelenumi{(\alph{enumi})}
\setcounter{enumi}{1}
\item
  Suppose now that \(\mathbf{K} = \mathbf{Z}_{\mathrm{P}}\) . Let \(c\) be an integer \(\geqslant 1\) . The group \(\mathbf{F}^c = \mathbf{F} / \mathbf{C}^c\mathbf{F}\) is identified with a subgroup of \(\Gamma /\Gamma_c\) by means of the homomorphism derived from \(g\) by taking quotients (notation of Theorem 2). The group \(\Gamma / \Gamma_c\) is a P-torsion-free P-divisible group of class \(< c\) . Let \(F_P^c\) be the P-saturation of \(F^c\) in \(\Gamma / \Gamma_c\) (§ 4, Exercise 14) and let \(i\) be the injection of \(F^c\) into \(F_P^c\) . The ordered pair \((i, F_P^c)\) is a P-envelope of \(F^c\) (§ 4, Exercise 15). Deduce that every nilpotent group has a P-envelope (note that every nilpotent group of class \(< c\) is the quotient of a group \(F^c\) for suitable X and use part (b) of Exercise 15 of § 4).
\item
  If \(n \leqslant c\) , let \(\mathbf{F}_{\mathbf{P},n}^{\mathrm{c}}\) be the intersection of \(\mathbf{F}_{\mathbf{P}}^{\mathrm{c}}\) with \(\Gamma_{n}/\Gamma_{c}\) ; if \(n \geqslant c\) , we write \(\mathbf{F}_{\mathbf{P},n}^{\mathrm{c}} = \{e\}\) . Show that \(\mathbf{F}_{\mathbf{P},n}^{\mathrm{c}}\) is the P-saturation of \(\mathbf{C}^{\mathrm{c}}\mathbf{F}^{\mathrm{c}}\) in \(\Gamma/\Gamma_{c}\) . The filtration \((\mathbf{F}_{\mathbf{P},n}^{\mathrm{c}})\) is an integral central filtration of \(\mathbf{F}_{\mathbf{P},n}^{\mathrm{c}}\) ; let \(\text{gr}(\mathbf{F}_{\mathbf{P}}^{\mathrm{c}})\) be the graded module associated with this filtration. Show that, if \(n < c\) , the image of \(\text{gr}_{\mathbf{n}}(\mathbf{F}_{\mathbf{P}}^{\mathrm{c}})\) in \(\text{gr}_{\mathbf{n}}(\Gamma) = \mathbf{A}_{\mathbf{Z}_{\mathbf{P}}}^{n}(\mathbf{X})\) is \(\mathbf{S}_{\mathbf{P}}^{-1}. \mathbf{L}_{\mathbf{Z}}^{n}(\mathbf{X}) = \mathbf{L}_{\mathbf{Z}_{\mathbf{P}}}^{n}(\mathbf{X})\) . Deduce that the Lie algebra \(\text{gr}(\mathbf{F}_{\mathbf{P},n}^{\mathrm{c}})\) is generated by its elements of degree 1 and hence that \(\mathbf{F}_{\mathbf{P},n} = \mathbf{C}^{n}(\mathbf{F}_{\mathbf{P}}^{\mathrm{c}})\) for all \(n\) ( \(\S 4\) , Exercise 1). Show that the group \(\mathbf{F}_{\mathbf{P}}^{\mathrm{c}}\) is generated by the \(x^{1/s}\) for \(x \in \mathbf{X}\) , \(s \in \mathbb{S}_{\mathbf{P}}\) (note that the images of these elements in \(\text{gr}_1(\mathbf{F}_{\mathbf{P}}^{\mathrm{c}}) \simeq \mathbf{Z}_{\mathbf{P}}^{(X)}\) generate the group \(\text{gr}_1(\mathbf{F}_{\mathbf{P}}^{\mathrm{c}})\) and apply Corollary 3 to Proposition 8 of \(Algebra\) , Chapter I, § 6, no. 3).
\item
  Let \(\mathrm{H}\) be a Hall set relative to \(\mathbf{X}\) (§2, no. 10) and let \(\mathrm{H}(c)\) be the subset of \(\mathrm{H}\) consisting of the elements of length \(< c\) . For \(m \in \mathrm{H}(c)\) , let \(\phi_c(m)\) denote the image in \(\mathbf{F}^c\) of the basic commutator \(\phi(m)\) defined by \(m\) (cf.~no. 4, Remark). Show that, for all \(w \in \mathbf{F}_P^c\) , there exists a unique element \(\alpha \in \mathbf{Z}_{\mathbf{P}}^{\mathbf{H}(c)}\) such that
\end{enumerate}

\[
w = \prod_ {m \in H (c)} \phi_ {c} (m) ^ {\alpha (m)}.
\]

(Use the determination of gr(F \(_{P}^{c}\) ) obtained above.)

\begin{enumerate}
\def\labelenumi{\arabic{enumi}.}
\setcounter{enumi}{6}
\tightlist
\item
  We preserve the notation of Exercise 6.
\end{enumerate}

\begin{enumerate}
\def\labelenumi{(\alph{enumi})}
\item
  Let G be a nilpotent group of class \(<c\) and let \((i, \overline{G})\) be a P-envelope of G. Let \(\rho: \mathrm{F}^{c} \to \mathrm{G}\) be a surjective homomorphism (such a homomorphism exists if X is suitably chosen) and let \(\bar{\rho}\) be the corresponding homomorphism of \(\mathrm{F}_{\mathbf{P}}^{c} \in \overline{G}\) (§ 4, Exercise 15). The homomorphism \(\bar{\rho}\) is surjective and maps \(\mathrm{C}^{n}\mathrm{F}_{\mathbf{P}}^{c} = \mathrm{F}_{\mathbf{P}, n}^{c}\) onto \(\mathrm{C}^{n}\overline{G}\) . Show that \(\mathrm{C}^{n}\overline{G}\) is the P-envelope of \(i(\mathrm{C}^{n}\mathrm{G})\) in \(\overline{G}\) and that the Lie algebra \(\mathrm{gr}(\overline{G})\) is identified with \(\mathbf{Z}_{\mathbf{P}} \otimes \mathrm{gr}(\mathrm{G})\) .
\item
  Deduce that G is P-divisible and P-torsion-free if and only if the \(C^{n}G/C^{n+1}G\) are.
\end{enumerate}

¶ 8. Suppose that K = Z and that X is finite. For every integer \(k \geqslant 0\) , let \((e_{k,\alpha})\) denote the basis of \(A_k(X)\) .

\begin{enumerate}
\def\labelenumi{(\alph{enumi})}
\tightlist
\item
  Let \((w_{n})_{n\in \mathbf{N}}\) be a sequence of elements of F. Let \(w_{k,\alpha}(n)\) be the coefficient of \(e_{k,\alpha}\) in the term of \(g(w_{n})\in \hat{\mathrm{A}} (\mathrm{X})\) of degree \(k\) . Then:
\end{enumerate}

\[
g (w _ {n}) = \sum_ {k, \alpha} w _ {k, \alpha} (n) e _ {k, \alpha} \quad \text { for   all } n \in \mathbf {Z}.
\]

The sequence \((w_{n})\) is called typical if, for every ordered pair \((k,\alpha)\) the function \(w_{k,\alpha}\colon n\mapsto w_{k,\alpha}(n)\) is a binomial polynomial of degree \(\leqslant k\) , cf.~Exercise 4. This condition is independent of the choice of the bases \((e_{k,\alpha})\) . Show that it is equivalent to the existence of a sequence \((a_{k})_{k\in \mathbf{N}}\) of elements of \(\hat{\mathbf{A}}_{\mathbf{q}}(\mathbf{X})\) such that \(\omega (a_k)\geqslant k\) for all \(k\) and

\[
g (w _ {n}) = \sum_ {k = 0} ^ {\infty} n ^ {k} a _ {k} \quad \text { for   all } n \in \mathbf {Z}.
\]

\begin{enumerate}
\def\labelenumi{(\alph{enumi})}
\setcounter{enumi}{1}
\item
  Show that, if \((w_{n})\) and \((w_{n}^{\prime})\) are two typical sequences, so are \((w_{n}^{-1})\) and \((w_{n}w_{n}^{\prime})\) .
\item
  Let \(w\) be an element of \(\mathbf{C}^k\mathbf{F}\) for \(k \geqslant 1\) and let \(f\) be a binomial polynomial of degree \(\leqslant k\) . Show that \((w^{f(n)})\) is a typical sequence. In particular if \(z \in \mathbf{F}\) , the sequence \((z^n)\) of powers of \(z\) is typical.
\item
  Let \((w_{n})\) be a sequence of elements of F. Show that \((w_{n})\) is typical if and only if there exist \(y_{0}, y_{1}, \ldots, y_{n}, \ldots\) in F, where \(y_{i} \in \mathbf{C}^{i}\mathbf{F}\) for all \(i \geqslant 1\) and
\end{enumerate}

\[
w _ {n} \equiv y _ {0} y _ {1} ^ {n} \dots y _ {k} ^ {\binom {n} {k}} \mod \mathrm{C} ^ {k + 1} \mathrm{F}
\]

for all \(n \in \mathbf{Z}\) and all \(k \geqslant 1\) . (Determine the \(y_{i}\) from \((w_{n})\) by setting successively \(n = 0, 1, \ldots\) . For example \(w_{0} = y_{0}, w_{1} = y_{0}y_{1}, \ldots\) )

\begin{enumerate}
\def\labelenumi{(\alph{enumi})}
\setcounter{enumi}{4}
\tightlist
\item
  Let H be a Hall set relative to X and let \((w_{n})\) be a sequence of elements of F. Let \((\alpha(m,n))_{m\in\mathbb{H},n\in\mathbf{Z}}\) be the family of integers such that
\end{enumerate}

\[
w _ {n} \equiv \prod_ {m \in \mathrm{H}} \phi (m) ^ {\alpha (m, n)} \mod \mathrm{C} ^ {c} \mathrm{F}
\]

for all \(n \in \mathbf{Z}\) and all \(c \geqslant 1\) (cf.~no. 4, Remark). Show that \((w_n)\) is typical if and only if, for all \(m \in \mathbf{H}\) , the function \(n \mapsto \alpha(m, n)\) is a binomial polynomial of degree \(\leqslant l(m)\) , where \(l(m)\) is the length of \(m\) .†

\begin{enumerate}
\def\labelenumi{(\alph{enumi})}
\setcounter{enumi}{5}
\tightlist
\item
  A sequence \((w_{n})\) is called 1-typical if the corresponding functions \(w_{k,\alpha}\) are binomial polynomials of degree \(\leqslant k - 1\) . Show that, if \((w_{n})\) is 1-typical, there exist \(y_{i} \in \mathbf{C}^{i + 1}\mathbf{F}\) for \(i = 0, 1, \ldots\) , such that
\end{enumerate}

\[
w _ {n} \equiv y _ {0} y _ {1} ^ {n} \dots y _ {k} ^ {\binom {n} {k}} \mod \mathrm{C} ^ {k + 1} \mathrm{F}
\]

for all \(n \in Z\) and all \(k \geqslant 1\) .

¶ 9. Let X be the set with two elements \{x, y\}.

\begin{enumerate}
\def\labelenumi{(\alph{enumi})}
\tightlist
\item
  Show that there exists a sequence \(w_{2}, w_{3}, \ldots\) of elements of \(\mathbf{F}\) such that \(w_{i} \in \mathbb{C}^{i}\mathbb{F}\) for all \(i\) and
\end{enumerate}

\[
(x y) ^ {n} \equiv x ^ {n} y ^ {n} w _ {2} ^ {\binom {n} {2}} \dots w _ {i} ^ {\binom {n} {i}} \mod \mathrm{C} ^ {i + 1} \mathrm{F}
\]

for all \(n \in \mathbf{Z}\) and all \(i \geqslant 1\) . (Apply Exercise 8 ( \(d\) ) to the sequence of \(x^{-n}(xy)^n\) .) Show that \(w_2 = y^{-1}(y^{-1}, x^{-1})y \equiv (x, y)^{-1} \mod \mathrm{C}^3\mathrm{F}\) .

\begin{enumerate}
\def\labelenumi{(\alph{enumi})}
\setcounter{enumi}{1}
\tightlist
\item
  Let G be a nilpotent group of class \(\leqslant c\) and let x, y be in G. Deduce from the above the following formula (``Hall's formula'');
\end{enumerate}

\[
(x y) ^ {n} = x ^ {n} y ^ {n} \prod_ {i = 2} ^ {c} w _ {i} (x, y) ^ {\binom {n} {i}}.
\]

\begin{enumerate}
\def\labelenumi{(\alph{enumi})}
\setcounter{enumi}{2}
\tightlist
\item
  Let \(p\) be a prime number. Show that there exists \(u_{i} \in \mathbf{C}^{i}\mathbf{F}\) ( \(2 \leqslant i \leqslant p - 1\) ) and \(w \in \mathbf{C}^{p}\mathbf{F}\) such that
\end{enumerate}

\[
(x y) ^ {p} = x ^ {p} y ^ {p} u _ {2} ^ {p} \dots u _ {p - 1} ^ {p} w.
\]

(Use (a), noting that the \(\binom{p}{i}\) is divisible by \(p\) if \(1 < i < p\) .) Let \(\bar{w}\) be the image of \(w\) in \(\mathrm{gr}_p(\mathsf{F}) = \mathrm{L}_{\mathbf{Z}}^p (\mathbf{X})\) and let \(\bar{w}_p\) be the image of \(\bar{w}\) in \(\mathrm{L}_{\mathbf{F}_p}^p (\mathbf{X})\) . Show that \(\bar{w}_p = \Lambda_p(x,y)\) , cf.~Chapter I, § 1, Exercise 19. (Use the extension \(g\) of F to \(\tilde{\mathbf{A}}(\mathbf{X})\) and compare the terms of degree \(p\) in \(g((xy)^p)\) and \(g(x^p y^p u_2^p \ldots u_{p-1}^p w)\) ). The first is equal to \((x + y)^p\) and the second is congruent mod. \(p\) to \(x^p + y^p + \bar{w}_p\) . Hence the result.)

Show that there exist \(a \in \mathbf{C}^2\mathbf{F}\) and \(w' \in \mathbf{C}^p\mathbf{F}\) such that

\[
(x y a) ^ {p} = x ^ {p} y ^ {p} w ^ {\prime}.
\]

(Use the above formula for \((xy)^p\) .) Deduce that in a nilpotent group of class \(p\) the \(p\) -th powers form a subgroup.

\begin{enumerate}
\def\labelenumi{(\alph{enumi})}
\setcounter{enumi}{3}
\tightlist
\item
  Show that there exists a sequence \(v_{3}, v_{4}, \ldots\) of elements of \(\mathbf{F}\) , where \(v_{i} \in \mathbf{C}^{i}\mathbf{F}\) for all \(i\) and
\end{enumerate}

\[
(x ^ {n}, y) \equiv (x, y) ^ {n} v _ {3} ^ {\binom {n} {2}} \dots v _ {i + 1} ^ {\binom {n} {i}} \mod \mathrm{C} ^ {i + 2} \mathrm{F}
\]

for all \(n \in \mathbf{N}\) and all \(i \geqslant 1\) . (Apply Exercise 8 (f) to the sequence of \((x^n, y)\) .)

Deduce that, if \(p\) is a prime number, there exists \(t_i \in \mathrm{C}^i\mathrm{F}\) for \(3 \leqslant i \leqslant p\) and \(z \in \mathrm{C}^{p+1}\mathrm{F}\) such that

\[
(x ^ {p}, y) = (x, y) ^ {p} t _ {3} ^ {p} \dots t _ {p} ^ {p} z.
\]

Show that the image \(\bar{z}_p\) of \(z\) in \(\mathrm{L}_{\mathbf{F}_p}^{p + 1}(\mathbf{X})\) is (ad \(x)^p (y)\) . (Same method as for \((c).)\)

¶ 10. Let p be a prime number and let G be a group with a real central filtration (G \(_{\alpha}\) ). Suppose that the relation x ∈ G \(_{\alpha}\) implies x \(^{p}\) ∈ G \(_{p\alpha}\) , in which case the filtration (G \(_{\alpha}\) ) is called restricted.

\begin{enumerate}
\def\labelenumi{(\alph{enumi})}
\item
  Show that the Lie algebra \(\operatorname{gr}(\mathbf{G})\) associated with \((\mathbf{G}_{\alpha})\) is such that \(p.\operatorname{gr}(\mathbf{G}) = 0\) and hence can be given an algebra structure over \(\mathbf{F}_p\) .
\item
  Let \(\xi \in \mathrm{gr}_{\alpha}(\mathbf{G})\) and let \(x\) be a representative of \(\xi\) in \(\mathbf{G}_{\alpha}\) . Show that the image of \(x^{p}\) in \(\mathrm{gr}_{p\alpha}(\mathbf{G})\) does not depend on the choice of \(x\) (use Exercise 9 (c)). If this image is denoted by \(\xi^{[p]}\) , prove that \(\xi \mapsto \xi^{[p]}\) is a linear mapping of \(\mathrm{gr}_{\alpha}(\mathbf{G})\) into \(\mathrm{gr}_{p\alpha}(\mathbf{G}')\) and that \((\xi + \xi')^{[p]} = \xi^{[p]} + \xi'^{[p]} + \Delta_p(\xi, \xi')\) . (Same method.)
\end{enumerate}

Show that, if \(\xi \in \mathrm{gr}_{\alpha}(\mathbf{G})\) and \(\eta \in \mathrm{gr}_{\beta}(\mathbf{G})\) , then

\[
[ \xi^ {[ p ]}, \eta ] = (\operatorname{ad} \xi) ^ {p} (\eta).
\]

(Use Exercise 9 (d).)

\begin{enumerate}
\def\labelenumi{(\alph{enumi})}
\setcounter{enumi}{2}
\tightlist
\item
  Show that there exists one and only one \(p\) -mapping of \(\operatorname{gr}(\mathbf{G})\) into itself (Chapter I, § 1, Exercise 20) which extends the mappings \(\xi \mapsto \xi^{[p]}\) defined on the \(\operatorname{gr}_{\alpha}(\mathbf{G})\) . With this structure, \(\operatorname{gr}(\mathbf{G})\) is a Lie \(p\) -algebra, said to be associated with the restricted filtration \((\mathbf{G}_{\alpha})\) .
\end{enumerate}

¶ 11. (a) Let A be a filtered algebra satisfying the conditions of §4, no. 5 and let \(\Gamma = A^{*} \cap (1 + A_{0}^{+})\) . \(\Gamma\) is given the filtration induced by that on A (loc. cit., Proposition 2). Show that, if K is a field of characteristic \(p > 0\) , \((\Gamma_{\alpha})\) is a restricted filtration (Exercise 10) and the embedding of \(\mathrm{gr}(\Gamma)\) in \(\mathrm{gr(A)}\) defined in loc. cit., Proposition 3 is compatible with the Lie \(p\) -algebra structures on \(\mathrm{gr}(\Gamma)\) and \(\mathrm{gr(A)}\) .

\begin{enumerate}
\def\labelenumi{(\alph{enumi})}
\setcounter{enumi}{1}
\item
  Suppose that \(\mathbf{K} = \mathbf{F}_p\) . We apply the above to the filtered algebra \(\hat{\mathbf{A}}(\mathbf{K})\) , in which \(\mathrm{F} = \mathrm{F}(\mathrm{X})\) is embedded by means of the homomorphism \(g\) . The filtration \((\Gamma_n)\) of \(\Gamma\) induces a filtration \((\mathbb{F}_n)\) on \(\mathbf{F}\) , which is a restricted filtration. The Lie \(p\) -algebra \(\operatorname{gr}(\mathbf{F})\) is identified with a Lie \(p\) -subalgebra of \(\mathrm{A}(\mathbf{X})\) containing \(\mathbf{X}\) . Deduce that \(\operatorname{gr}(\mathbf{F})\) contains the free Lie \(p\) -algebra \(\mathrm{L}(\mathbf{X}, p)\) , cf.~§ 3, Exercise 4 (e).
\item
  Let \(\mathbf{H}\) be a Hall set relative to \(\mathbf{X}\) ; for every integer \(i\) , let \(\mathbf{H}_i\) be the set of elements of \(\mathbf{H}\) of length \(i\) . Let \(w \in \mathbf{F}\) and let \(\alpha_i \in \mathbf{Z}^{(\mathbf{H}_i)}\) be such that
\end{enumerate}

\[
w \equiv \prod_ {i = 1} ^ {n} \prod_ {m \in H_i} \phi (m) ^ {\alpha_ {i} (m)} \mod C ^ {n + 1} F
\]

for all \(n\) , cf.~no. 4. For all \(m \in \mathbf{H}\) , let \(l(m)\) denote the length of \(m\) and \(q(m, w)\) the largest power of \(p\) which divides \(\alpha_{1}(m)\) , where \(i = l(w)\) (if \(\alpha_{1}(m) = 0\) , we make the convention \(q(m, w) = +\infty\) ). Let \(\mathrm{N} = \inf_{m \in \mathbb{H}} (l(m).q(m, w))\) . Show that the \(\phi(m)^{\alpha_{1}(m)}\) belong to \(\mathrm{F}_{\mathrm{N}}\) and even to \(\mathrm{F}_{\mathrm{N+1}}\) if \(l(w).q(m, w) > \mathrm{N}\) . Show that, if \(l(w).q(m, w) = \mathrm{N}\) , the image of \(\phi(m)^{\alpha_{1}(m)}\) in \(\mathrm{gr}_{\mathrm{N}}(\hat{\mathrm{A}}(\mathbf{X})) = \mathrm{A}^{\mathrm{N}}(\mathbf{X})\) is a nonzero multiple of \(\overline{m}^{q(m, w)}\) , where \(\overline{m}\) is the element of \(\mathrm{L}(\mathbf{X})\) defined by \(m\) ( \(\S 2\) , no. 11). Now these elements are linearly independent ( \(\S 3\) , Exercise 4); deduce that the image of \(w\) in \(\mathrm{gr}_{\mathrm{N}}(\mathbf{F})\) is \(\neq 0\) and belongs to \(\mathrm{L}(\mathbf{X}, p)\) , whence the fact that \(\mathrm{gr}(\mathbf{F}) = \mathrm{L}(\mathbf{X}, p)\) .

\begin{enumerate}
\def\labelenumi{\arabic{enumi}.}
\setcounter{enumi}{11}
\tightlist
\item
  Let g be a Lie algebra over a field k of characteristic p \textgreater{} 0.
\end{enumerate}

\begin{enumerate}
\def\labelenumi{(\alph{enumi})}
\tightlist
\item
  Let \(h\) be an integer \(\leqslant p - 1\) . \(\mathfrak{g}\) is said to satisfy Engel's \(h\) -th condition if \((\operatorname{ad} x)^h(y) = 0\) for every ordered pair \((x, y)\) of elements of \(\mathfrak{g}\) . Show that this is equivalent to
\end{enumerate}

\[
\sum_ {\sigma \in \mathfrak {S} _ {h}} \operatorname{ad} (x _ {\sigma (1)}) \circ \dots \circ \operatorname{ad} (x _ {\sigma (h)}) = 0
\]

for \(x_1, \ldots, x_h\) in \(\mathfrak{g}\) (apply the formula \((\operatorname{ad} x)^h = 0\) to the linear combinations of the \(x_i\) ).


\begin{enumerate}
\def\labelenumi{(\alph{enumi})}
\setcounter{enumi}{1}
\item
  Let \(x, y\) be elements of \(\mathfrak{g}\) such that \(\Lambda_p(ax, by) = 0\) for all \(a, b\) in \(k\) . Show that, if \(\Lambda_p^{r,s}\) is the bihomogeneous component of \(\Lambda_p\) of bidegree \((r, s)\) , then \(\Lambda_p^{r,s}(x, y) = 0\) for every ordered pair \((r, s)\) . Deduce (taking \(r = p - 1\) , \(s = 1\) ) that \((\text{ad } x)^{p-1}(y) = 0\) . In particular, if \(\Lambda_p(x, y) = 0\) for \(x, y\) in \(\mathfrak{g}\) , the Lie algebra \(\mathfrak{g}\) satisfies Engel's \((p - 1)\) -th condition.
\item
  Let \(\mathfrak{c}\) be the centre of \(\mathfrak{g}\) . Suppose that \((\operatorname{ad} x)^p = 0\) for all \(x \in \mathfrak{g}\) . Show that \(\mathfrak{g} / \mathfrak{c}\) satisfies Engel's \((p - 1)\) -th condition. (Show that \(\Lambda_p(\operatorname{ad} x, \operatorname{ad} y) = 0\) for \(x, y\) in \(\mathfrak{g}\) and apply \((b)\) to the Lie algebra \(\operatorname{ad} \mathfrak{g} = \mathfrak{g} / c\) .)
\item
  Let G be a group such that \(x^{p}=1\) for all \(x\in G\) and let \((\mathrm{G}_{\alpha})\) be a real central filtration of G. The filtration \((\mathrm{G}_{\alpha})\) is restricted (Exercise 10) and \(\mathrm{gr}(\mathrm{G})\) is a Lie p-algebra whose p-mapping is zero. Deduce that \(\mathrm{gr}(\mathrm{G})\) satisfies Engel's \((p-1)\) -th condition.†
\end{enumerate}

